% Paper 2, Section 8 — Hardware implications (evidence E6 + the energy columns of E1–E3).
% Numbers verified against ledger entries 83e110a (ADC-step-floors-theta), a9f36d2 (usable
% theta window on copy; theta=0.5 collapse), 7a10a0d (theta non-transfer across tasks),
% bb79f7f (char-LM energy/quality tradeoff vs regularized baseline), a8e7a7c / 8e74f48 /
% 159e341 (copy energy columns), 5bb7bcc / 8d050cb / 48510c8 (margins vs regularized reference),
% c8b596f / a9ff715 (energy decomposition, free-recurrence ceiling, converter/storage sweep),
% 73c99e6 (copy event-readout refutation), 60d627c (outrand attribution), db7c13b (3-seed
% confirmation of all gated cells — the n=1 caveat is retired for the cited claims),
% d47511b (shape law sweep), 3033631/c18c724 (streaming row, pre-registered REFUTE),
% 39a6873/0924aad (budget check, pre-registered CLOSE — deficit flat ep10->ep30).
% CLAIMS DISCIPLINE — forbidden in this section, per the skeleton's claims ledger:
%   "analog beats digital" (retracted); any pJ number presented as a silicon measurement
%   (all energy figures are the 45nm Horowitz MAC/AC proxy); quoting copy margins vs the
%   PLAIN digital reference (corrected margins vs the noise-regularized reference only).

\section{Hardware implications}
\label{sec:hardware}

The preceding sections were organized around a scientific question --- which datapath
degradations does a workload tolerate? This section reads the same results as design
guidance: which route should carry a state-space model on neuromorphic silicon, and what
co-design constraints bind once a route is chosen.

\subsection{The recommendation splits by workload}
\label{sec:hardware-split}

No route wins both tasks, and the ordering inversion of Section~\ref{sec:copy} is stable at
three seeds on both sides, so the recommendation cannot be collapsed to a single answer.

\emph{Statistical sequence workloads with low retention demand} tolerate a lossy recurrent
state where we measured it on char-LM: the analog route holds the state in analog dynamics,
communicates on 61\% of steps, and pays $+0.160$ bpc (5.3\%) against a noise-regularized
digital baseline for a $1.60\times$ proxy-energy reduction
(Section~\ref{sec:charlm-claim}). This is a tradeoff, not a free lunch, and it is the only
cell in any task where a neuromorphic datapath buys energy at single-digit quality cost ---
the tolerance is not generic to statistical workloads: on a streaming perception task at a
small-output shape the same analog state costs $10$--$19$ accuracy points, budget-robustly
(Section~\ref{sec:hardware-shape}).

\emph{Precise-recall workloads} invert the choice. On copy the analog route loses to
output-spiking on \emph{both} axes (Section~\ref{sec:hardware-energy}), and the route that
survives is the one SPikE-SSM instantiates: keep the state in an exact
digital datapath and spike only the output, which retains 0.87 and 0.42 of the regularized
baseline's above-chance margin at $M{=}16$ and $M{=}32$ while emitting on $\sim$45\% of steps.

\emph{Carrying the recurrent state in spikes is a non-option for any memory-bearing
workload.} At $M{=}32$ and again at $M{=}64$ the spiking-state variant sits at exactly
chance ($0.0623\pm0.0012$ and $0.0623\pm0.0005$ against 0.0625) --- a 1-bit state does not
degrade precise recall, it eliminates it (Section~\ref{sec:copy-floor}). The live design
fork on real silicon is therefore binary: an exact digital state with a spiked output, or an
analog state that accepts the loss and gains the event sparsity. Which side of the fork is
right is set by the workload's retention demand, not by the hardware.

\subsection{Event sparsity and state precision are not independent knobs}
\label{sec:hardware-adc}

The analog route's operating point is set by the send-on-delta threshold $\theta$, and the
natural assumption is that $\theta$ is a free knob: lower it for quality, raise it for
sparsity. Two measured constraints break that assumption.

First, \emph{the ADC step floors the event rate}. The simulated analog state is quantized to
$b{=}6$ bits over rails $\pm r$ with $r{=}4$, giving a least significant bit of
$q = 2r/2^{b} = 0.125$, and the send-on-delta comparison is made on the \emph{quantized}
state. Every $\theta$ below $q$ is therefore inoperative: $\theta{=}0.05$ and
$\theta{=}0.10$ produce bit-identical results on every metric to all recorded digits. Below
one LSB the quantizer, not the threshold, decides which state changes are visible, so the
minimum achievable event rate is set by converter precision. Buying a lower event rate means
spending bits --- a converter area and energy cost --- not tuning a threshold. This couples
two quantities a designer would prefer to optimize separately, and it is invisible in any
formulation that treats the state as continuous.

Second, \emph{$\theta$ does not transfer across workloads}. The char-LM-calibrated
$\theta{=}1.0$ collapses the analog variant to chance on copy, emitting on only 3--8\% of
steps; the usable window on copy is $0.125 < \theta \lesssim 0.3$, and $\theta{=}0.5$
already over-gates to chance. The window's lower edge is the LSB constraint above; its upper
edge is task-dependent. An analog SSM deployment therefore needs per-workload threshold
calibration, and a fixed-$\theta$ part cannot serve both workload classes of
Section~\ref{sec:hardware-split}.

\subsection{What the energy proxy does and does not license}
\label{sec:hardware-energy}

All energy figures in this paper are the 45nm proxy of Section~\ref{sec:energy} (4.6 pJ per
MAC, 0.9 pJ per AC); none are silicon measurements. Read under that caveat, the proxy
licenses one claim and blocks two others.

It licenses the char-LM tradeoff: analog at $\theta{=}0.15$ costs 446k pJ/token against the
regularized digital baseline's 712k --- a $1.60\times$ reduction, purchased for $+0.160$
bpc. The digital baseline's noise regularizer is training-time only, so its inference
datapath and energy are identical to the plain baseline's; the comparison is fair on both
axes.

It blocks any analog energy claim on copy. At $M{=}16$ the output-spiking variant keeps
0.87 of the regularized margin at 122.9k pJ/token ($3.2\times$ below digital's 398.0k),
while analog keeps 0.50 at 246.7k --- worse on quality \emph{and} energy; at $M{=}32$ the
pattern repeats (0.42 at 125.2k versus 0.28 at 231.0k). The datapath reason is structural:
the analog variant earns event pricing only on the recurrent mixing matrix, while its input
and output projections remain MAC-priced, so its proxy energy cannot approach
output-spiking's unless most parameters sit in the recurrence.

And it blocks reading cheapness as merit. The cheapest cell in the entire copy table is the
spiking-state variant at $M{=}32$ (102k pJ/token) --- the same cell that retains exactly
nothing. A pJ/token figure is meaningless without the quality it purchased.

Finally, the proxy's base accounting omits costs that fall specifically on the analog
route: the analog state is dense (state density 0.994), and each unit carries an analog
storage element and a converter that MAC/AC counting alone does not price. Matching
variants on emitted rate matches wire traffic, not state-holding cost
(Section~\ref{sec:metrics}). Section~\ref{sec:hardware-stranded} prices both terms under
the proxy and finds them negligible at this width --- but that is a proxy statement, not a
silicon one, which is why measured Loihi-2/SpiNNaker-2-class energy, not more simulation,
remains the natural next step.

\subsection{The energy is not where the mechanisms are}
\label{sec:hardware-stranded}

The structural observation above --- that the analog route earns event pricing only on the
recurrent mixing matrix --- invites an obvious response: make the rest of the datapath
event-driven too. We attempted exactly that, and the attempt failed in a way that is
itself a design finding, so we report it in full. All gated-cell effects below are
three-seed means, paired against each arm's own no-gate reference.

\emph{Decomposing the proxy locates the energy away from every mechanism in this paper.}
On char-LM, the readout projection $W_{\mathrm{out}}$ carries 75\% of the analog variant's
proxy energy per token (and 47\% of digital's), while the recurrence $W_{\mathrm{mix}}$
--- the only term the proxy prices as event-driven for analog --- carries 8\%. The reason
is dimensional: at vocabulary $V{=}284$ the readout ($H \times V = 72{,}448$ weights) is
the largest single matrix in the model, exceeding the recurrence ($H^{2} = 65{,}536$).
The consequence is a ceiling: an analog variant with a \emph{free} recurrence (event rate
$\to 0$) still costs 410{,}040 pJ/token, a $1.74\times$ reduction against the regularized
digital reference --- so the published $1.60\times$ of
Section~\ref{sec:hardware-energy} already realizes 92\% of everything recurrence
sparsification can ever deliver on this datapath. Every neuromorphic mechanism examined in
this paper, and every route surveyed in Section~\ref{sec:related}, acts on the 8\% term.

\emph{Converters are not the obstacle.} Extending the proxy to the terms it previously
omitted --- ADC/DAC conversion swept over 0.01--1.0 pJ per graded event, the analog
storage element over 0--0.5 pJ per unit per step --- leaves both at 0.00--0.33\% of the
analog total at $H{=}256$, two to four orders of magnitude below the matrix terms. This
cuts both ways: it removes unpriced converters as an objection to the analog route at this
width, and it does not license calling converters free in general --- they begin to matter
exactly when the matrices become cheap.

\emph{The readout's share is stranded: an event-driven readout is quality-fatal on both
workload classes.} We implemented an incremental send-on-delta readout --- each unit
transmits into $W_{\mathrm{out}}$ only when its value has moved by more than a threshold
since last sent (verified exact at a near-zero threshold) --- and swept it on both tasks.
On char-LM the gate costs $+0.99$ to $+1.12$ bpc on the digital arm and $+0.91$ to $+1.05$
on the analog arm at \emph{every} threshold, flat across send rates $0.98 \to 0.11$: a
step function at gate-on, not a price per unit of sparsity bought. The digital arm pays it
equally, so none of the readout saving is attributable to the analog state. On copy the
same gate does not degrade recall, it eliminates it --- margin kept $\leq 0.05$ at both
$M{=}16$ and $M{=}32$ at every threshold, including one where 99.8\% of readout units
still transmit; the gated cells train to the uniform distribution or diverge, so the gate
breaks training, not merely inference.

\emph{Staleness itself, not the selection rule, is the fatal mechanism.} Holding readout
units by a random per-step Bernoulli draw at matched send rates costs $+1.12$/$+0.92$ bpc
(digital/analog) already at a 0.90 send rate, and the delta-triggered rule is
\emph{better} than random at low rates (by $0.45$ bpc at a $0.08$ send rate, digital arm)
--- send-on-delta is exonerated as a selection rule, and no smarter gating rule reopens
the route. The one readout sparsification that survives quality --- LIF output-spiking,
which keeps 0.87/0.42 of the regularized margin on copy
(Section~\ref{sec:hardware-split}) --- recomputes the full readout every step, and
therefore recovers none of its MAC energy.

The design conclusion is blunt. On a language-like workload the pJ ceiling of this SSM
datapath family is set by the readout width, not by the state mechanism, and the readout's
75\% share resists both known gating mechanisms: hold-based schemes are quality-fatal, and
spike-based schemes save no arithmetic. Claims of large energy wins from sparsifying an
SSM recurrence should be read against the 8\% share that term actually carries \emph{at a
language-model shape} --- Section~\ref{sec:hardware-shape} shows that both the share and
the ceiling are set by the readout-to-state ratio $V/H$, and tests the shape where the
projection favours the recurrence.

\subsection{The ceiling is a shape property --- and the shape that recovers the energy
fails on quality}
\label{sec:hardware-shape}

The $1.74\times$ ceiling above is a property of $V/H = 1.11$, not of the analog datapath.
Sweeping the same accounting (identical terms, nothing re-derived) over readout width $V$
at fixed $E{=}64$, $H{=}256$ and the measured char-LM operating point ($r_z = 0.6122$,
three seeds): at $V{=}10$ ($V/H = 0.04$) the readout falls to 9.5\% of the analog total,
the recurrence becomes the \emph{largest} single term (29.2\%), and the same measured
activity is worth a projected $3.16\times$ (ceiling $4.46\times$); at $V{=}284$ (char-LM),
$1.60\times$ (ceiling $1.74\times$); at $V{=}1024$, $1.20\times$; at a word-piece
vocabulary ($V{=}32{,}000$), $1.01\times$ --- nothing. Two corollaries sharpen the sweep.
The published $1.60\times$ realizes 92\% of its own shape's ceiling, so there was never
anything left to win at char-LM, whereas at $V{=}10$ the measured activity realizes only
71\% of the ceiling. And at $V{=}10$ the projected advantage \emph{grows} with $H$
($1.70\times$/$3.16\times$/$5.58\times$ at $H{=}64$/$256$/$1024$), because the
event-priced term is the $H^{2}$ one --- the opposite of the char-LM regime. Pricing the
input as an event stream (symmetrically: it is equally available to the digital arm) is
worth $+0.12\times$ at $V{=}284$ but $+3.06\times$ at $V{=}10$, where removing the readout
blocker promotes $W_{\mathrm{in}}$ to 61\% of the analog total; the two effects are
multiplicative, reaching $6.21\times$ at $V{=}10$ with input event rate $0.2$. Converters
re-checked at their most favourable corner ($V{=}10$ with event input) reach 0.58\% ---
still not decisive, now as a bounded statement.

The projection therefore locates the one shape where a large pJ win is available, and no
quality number had ever been measured there. We measured it, with criteria pre-registered
before any cell ran: FashionMNIST row streaming (28 steps of 28 pixels, $V{=}10$,
$V/H = 0.039$, chance accuracy $0.10$, parameter-matched at $86{,}858$). The
\emph{energy} side of the shape law confirms out of sample: measured $3.29\times$ against
the projected $3.16\times$, and the digital arm's decomposition inverts exactly as
predicted ($W_{\mathrm{mix}}$ 75.7\%, $W_{\mathrm{in}}$ 21.0\%, $W_{\mathrm{out}}$ 3.0\%).
The \emph{quality} side is refuted: against the noise-regularized reference (accuracy
$0.7878$), the best analog setting ($\theta{=}0.15$, state event rate $0.428$) reaches
$0.6774$ --- an 11.0-point deficit, and the deficit exceeds the pre-registered 5-point
criterion at every threshold ($11.0$/$19.0$/$25.7$/$42.3$/$58.3$ points at
$\theta = 0.15$--$2.0$). Unlike the copy collapse this is a graded tradeoff, not a
training failure --- the event rate falls smoothly $0.43 \to 0.05$ --- but no point on
the curve is quality-matched.

The deficit is budget-robust. Because the copy row had taught us that budget-limited
baselines mislead, we re-ran the reference and both leading analog cells at three times
the epochs and the full training set ($30$ epochs, $60$k images, pre-registered): the
reference improves to $0.8442$ (the training-time noise benefit of
Section~\ref{sec:charlm-control} reappears at the long budget, reversing the short-budget
ordering), the analog cells improve in step, and the deficit is flat
($11.0 \to 10.4$ and $19.0 \to 18.9$ points). Tripling the budget moves the gap by less
than $0.6$ points, so the deficit is a property of the lossy analog state at this shape,
not of an under-trained reference.

Three readings, in decreasing order of importance. First, the $3.3$--$4.0\times$ band is
\emph{shape-determined, not mechanism-determined}: all three neuromorphic variants land in
it (output-spiking $3.34\times$ at accuracy $0.6486$, spiking-state $3.63\times$ at
$0.5298$), because at $V{=}10$ the recurrence dominates and every mechanism earns event
pricing on it --- no mechanism can claim the factor. Second, the analog route is the best
non-digital variant here (the char-LM ordering reproduces; the copy inversion does not,
consistent with Section~\ref{sec:hardware-split}), but only by 2.9 points over
output-spiking --- and output-spiking's $3.34\times$ carries no graded-event pricing
ambiguity, while analog's $3.29\times$ falls to $1.77\times$ if each graded event is
conservatively priced as a MAC. Priced conservatively, the defensible operating point at
this shape is a spiked output on an exact state, not an analog state. Third, the principle
of Section~\ref{sec:principle} predicted the opposite outcome here --- a state degradation
was expected to be cheap on a statistical workload --- which is its second failed forecast
(Section~\ref{sec:principle-against}).

Across every shape tested, the recoverable energy and the tolerable quality loss never
co-occur: at the language shape the quality cost is single-digit but the readout strands
the energy; at the streaming shape the energy is real and measured but the analog state
costs double-digit accuracy at every threshold and every budget tested. The verdict is
therefore unqualified within the tested space: \emph{no quality-matched pJ win at any
model shape tested, at any training budget tested.} Caveats: the streaming cells are
single-seed (the effects, $10$--$19$ points, dwarf the $\leq 1$-point seed noise
everywhere else in this pipeline); the input is dense, so the $6.21\times$ event-input
corner is untested --- but it improves only the energy axis, which is not the axis that
failed; and all figures remain the 45\,nm proxy.
